\SourcePage{10}{15}
\section{Uncertainty Relations in the Relativistic Domain}

The quantum theory presented in Vol.~III of this course is non-relativistic in its very substance, and it fails for phenomena in which motion occurs at speeds no longer small next to the speed of light. One might at first imagine that a relativistic theory could be reached by generalizing the formalism of non-relativistic quantum mechanics in a more or less straightforward manner. A closer look shows, however, that a logically complete relativistic theory cannot be built without invoking fresh physical principles.

We begin by recalling several physical notions on which non-relativistic quantum mechanics rests (Vol.~III, \S~1). A fundamental part was played there, as we saw, by the concept of measurement~--- the process in which a quantum system interacts with a ``classical object'' (an ``apparatus'') and thereby comes to possess definite values of particular dynamical variables (coordinates, velocities, and so forth). We saw, too, that quantum mechanics places severe limits on the extent to which different dynamical variables of an electron\footnote{As in Vol.~III, \S~1, we say ``electron'' for brevity, understanding by it any quantum system.} can exist together. Thus the uncertainties $\Delta q$ and $\Delta p$ with which coordinate and momentum can coexist obey the relation $\Delta q \Delta p \sim \hbar$;\footnote{Ordinary units are used in this section.} the more precisely either of these quantities is measured, the less precisely the other can be known at the same moment.

What matters, however, is that any single dynamical variable of the electron, on its own, admitted measurement of unlimited precision, and within a time interval that could be taken as short as desired. On this circumstance the whole of non-relativistic\SourcePage{11}{16} quantum mechanics depends in a fundamental way. It alone makes it possible to introduce the wave function, the central object of the theory's formalism. The physical content of the wave function $\psi(q)$ is, after all, that its squared modulus gives the probability that a measurement performed at a given instant will yield one or another value of the electron's coordinate. Evidently the very notion of such a probability presupposes that a coordinate measurement of unlimited precision and speed can be realized in principle; failing that, the notion would be empty of content and would forfeit its physical meaning.

The existence of a limiting speed (the speed of light~$c$) leads to new limitations, of a fundamental kind, on the possibilities of measuring various physical quantities (\emph{L.~D.~Landau, R.~Peierls}, 1930).

In Vol.~III, \S~44 the relation
\begin{equation}
(v' - v)\,\Delta p\,\Delta t \sim \hbar
\tag{1.1}
\end{equation}
was obtained; it relates the uncertainty $\Delta p$ of a momentum measurement on the electron to the length $\Delta t$ of the measuring process itself, $v$ and $v'$ denoting the electron's velocity before and after that process. It follows that momentum can be measured accurately within a short time (small $\Delta p$ together with small $\Delta t$) only if the measurement itself alters the velocity strongly enough. In the non-relativistic theory this expresses the fact that a momentum measurement repeated after a short interval will not reproduce its result, but it places no restriction of principle on how accurate a single measurement of the momentum may be, the difference $v' - v$ being capable of unlimited increase.

The presence of a limiting speed, however, changes the state of affairs radically. The difference $v' - v$, like the velocities themselves, can now not exceed $c$ (more precisely,~$2c$). Replacing $v' - v$ in~(1.1) by~$c$, we obtain the relation
\begin{equation}
\Delta p\,\Delta t \sim \hbar/c,
\tag{1.2}
\end{equation}
which determines the best accuracy of momentum measurement attainable in principle for a given duration $\Delta t$ of the measurement. Thus in the relativistic theory an arbitrarily accurate and arbitrarily rapid measurement of the momentum turns out to be impossible in principle. An exact measurement of the momentum ($\Delta p \to 0$) is possible only in the limit of an infinitely long duration of the measurement.

One has reason to expect that the measurability of the electron's coordinate as such is likewise affected. Within\SourcePage{12}{17} the mathematical formalism of the theory this appears as a conflict between an exact coordinate measurement and the statement that a free particle's energy is positive. As we shall see, the full set of eigenfunctions of a free particle's relativistic wave equation contains, besides the solutions whose time dependence is the ``proper'' one, solutions of ``negative frequency'' as well. These will in general appear also in the expansion of a wave packet describing an electron confined to a small region of space.

As will be shown, the ``negative-frequency'' wave functions are tied to the existence of antiparticles~--- the positrons. Their presence in the expansion of the wave packet reflects the creation of electron--positron pairs, unavoidable in the general case, during a measurement of the electron's coordinates. Since the emergence of these new particles cannot be controlled by the measuring process itself, coordinate measurement for the electron is drained of meaning.

In the rest frame of the electron the minimum error in a measurement of its coordinates is
\begin{equation}
\Delta q \sim \hbar/mc.
\tag{1.3}
\end{equation}
This value (the only one that dimensional considerations would permit) answers to a momentum uncertainty $\Delta p \sim mc$, which corresponds in turn to the minimum threshold energy needed to create a pair.

In a reference frame where the electron moves with energy~$\varepsilon$, (1.3) is replaced by
\begin{equation}
\Delta q \sim c\hbar/\varepsilon.
\tag{1.4}
\end{equation}

In the extreme ultra-relativistic limit, in particular, energy and momentum are related by $\varepsilon \approx cp$, so that
\begin{equation}
\Delta q \sim \hbar/p,
\tag{1.5}
\end{equation}
the error $\Delta q$ then coinciding with the particle's de~Broglie wavelength.\footnote{What is meant here are measurements any outcome of which permits an inference about the state of the electron; we thus set aside coordinate measurements by collision, in which over the time of observation the result occurs with a probability short of unity. In such a case a deflection of the particle does allow the electron's position to be inferred, but from the absence of a deflection nothing at all can be concluded.}

For photons the ultra-relativistic case always holds, so that expression~(1.5) is valid. This means that it makes sense to speak of the coordinates of a photon only in those cases where the characteristic dimensions are large compared with the wavelength. But this is nothing other than the ``classical'' limiting case, corresponding to geometrical optics, in which\SourcePage{13}{18} one may speak of the propagation of light along definite trajectories~--- rays. In the quantum case, however, when the wavelength cannot be regarded as small, the concept of the coordinates of a photon becomes vacuous. We shall see later (see~\secref{4}) that in the mathematical formalism of the theory the non-measurability of the photon's coordinates already shows itself in the impossibility of forming from its wave function a quantity that could play the part of a probability density satisfying the requirements which relativistic invariance demands.

On the basis of all that has been said, it is natural to suppose that the future theory will renounce altogether the consideration of the time development of the processes of interaction of particles. It will show that in these processes no exactly definable characteristics exist (even within the limits of the usual quantum-mechanical accuracy), so that the description of a process in time will turn out to be just as illusory as classical trajectories proved to be in non-relativistic quantum mechanics. The only observable quantities will be the characteristics (momenta, polarizations) of free particles~--- the initial particles that enter into the interaction and the final particles that have arisen as a result of the process (\emph{L.~D.~Landau, R.~Peierls}, 1930).

The characteristic formulation of a problem in relativistic quantum theory consists in the determination of the probability amplitudes of transitions connecting given initial and final states (that is, at $t \to \mp\infty$) of a system of particles. The totality of the transition amplitudes between all possible states constitutes the \emph{scattering matrix}, or \emph{$S$-matrix}. This matrix will be the carrier of all the information about the processes of interaction of particles that possesses an observable physical meaning (\emph{W.~Heisenberg}, 1938).

At the present time a complete, logically closed relativistic quantum theory does not yet exist. We shall see that the existing theory introduces new physical aspects into the manner of describing the states of particles, which acquires certain features of field theory (see~\secref{10}). The theory is constructed, however, to a considerable degree on the model of, and with the help of the concepts of, ordinary quantum mechanics. This way of building the theory has led to success in the domain of quantum electrodynamics. The absence of complete logical closure shows itself in this theory in the existence of divergent expressions when its mathematical apparatus is applied directly; but for the removal of these divergences there exist entirely unambiguous procedures. Nevertheless these procedures retain to a considerable degree the character of semi-empirical recipes, and our confidence in the correctness of the results obtained in this way rests, in the end, on their excellent agreement with experiment, and not on the internal consistency and logical coherence of the fundamental principles of the theory.
